\section{Threats to Validity}
\label{sec:threats}

\textbf{Internal validity.} The main internal risk is that the hosted model may change after our experiment. If OpenAI changes serving behavior behind a model alias, a later run may not reproduce the same feature files. We reduced that risk by using \texttt{gpt-5.4-mini-2026-03-17}, setting \texttt{temperature=0}, and keeping the API trace in \texttt{results/full\_api\_log.jsonl}. This is not complete control over the model provider, but it is enough to show which configuration produced our reported files. Another internal risk comes from the reverse-engineering step used to convert existing Gherkin into Connextra-style user stories. If that reverse story is too close to the original feature file, the task may become easier than writing Gherkin from a real backlog item. We reduced this risk by keeping the user story format focused on role, goal, benefit, and acceptance criteria rather than copying full Gherkin steps.

\textbf{External validity.} The dataset is multi-domain, not universal. Sylius, Apache Fineract, and Diaspora are public projects with mature Gherkin suites. Their scenarios are cleaner than many classroom or early-stage Jira tickets. Results may be lower for teams that write short free-form requirements, mix Vietnamese and English in one backlog, or lack a stable set of BDD step definitions. We therefore treat the result as evidence for public English-language BDD repositories rather than for every agile team. The domain balance is also imperfect: Sylius and Fineract contribute 100 cases each, while Diaspora contributes 61. This means the full median is influenced more by e-commerce and finance than by social-network stories.

\textbf{Construct validity.} Skeleton cosine similarity can miss important business differences. A generated Fineract scenario might keep the same Given--When--Then structure but change a repayment condition. A Sylius scenario might parse correctly while forgetting the channel-specific shipping constraint. To reduce this problem, we paired cosine with \texttt{behave --dry-run}. The parser check tells us whether the output is valid Gherkin syntax. It does not tell us whether the business rule is complete, so human review remains necessary. The skeleton preprocessing itself is another construct choice. Replacing quoted values and numbers makes the comparison less sensitive to surface data, but it may also hide errors where the exact date, price, or amount is the core of the scenario.

\textbf{Conclusion validity.} The statistical decision depends on the two thresholds chosen before the full run: 0.85 for skeleton cosine similarity and 80\% for executable syntax rate. RQ1 uses a one-sample Wilcoxon signed-rank test; RQ2 uses an exact binomial test. The full dataset contains 261 cases, so the result is not driven by the 30-sample pilot. The main limitation is the semantic cutoff itself. There is no accepted Gherkin benchmark saying that 0.85 cosine is the boundary for expert-level behavior. For that reason, we report both the median and the hypothesis-test result instead of treating the threshold as a final quality label.
